\chapter{Evaluation Framework}\label{ch:evaluation}

\section{Introduction}

Every result in this report comes from one framework, built so that a result can only be made one way and can always
be traced back to the raw records of its trials. This chapter describes it; Chapter~\ref{ch:design} describes the
experiments run with it.

\section{Three kinds of evaluation}\label{sec:three-evals}

Figure~\ref{fig:pipeline} places the three kinds of evaluation in the project: every trained policy is evaluated
offline, on the robot and on the replay bench.

\begin{figure}[htbp]
  \centering
  % The whole project: record -> train -> evaluate (three ways) -> report, and what carries the work between stages.
% Shared figure colours (colour only in figures).
\definecolor{robotfill}{HTML}{EAF2FC}\definecolor{robotline}{HTML}{2A78D6}
\definecolor{serverfill}{HTML}{FDEEE7}\definecolor{serverline}{HTML}{EB6834}
\definecolor{netfill}{HTML}{F0EFEC}\definecolor{greenfill}{HTML}{E9F5EC}\definecolor{greenline}{HTML}{2F9E58}
\definecolor{purplefill}{HTML}{F1ECF9}\definecolor{purpleline}{HTML}{7A55C7}
\definecolor{ink}{HTML}{0B0B0B}\definecolor{ink2}{HTML}{52514E}\definecolor{faultc}{HTML}{C9C6BF}

\begin{tikzpicture}[font=\small, node distance=5mm and 7mm,
  st/.style={draw=#1, thick, rounded corners=2pt, fill=white, align=center, minimum width=29mm, minimum height=13mm,
             inner sep=3pt},
  arr/.style={-{Latex[length=2.2mm]}, thick, ink2}]
\node[st=robotline] (rec) {\textbf{Record}\\ leader arm, cameras\\ trim, review};
\node[st=serverline, right=of rec] (tr) {\textbf{Train}\\ 7 policies,\\ original recipes};
\node[st=greenline, right=of tr, yshift=17mm] (off) {\textbf{Offline}\\ held-out error,\\ inference cost};
\node[st=greenline, right=of tr] (rob) {\textbf{Robot}\\ trials, success,\\ the conditions};
\node[st=greenline, right=of tr, yshift=-17mm] (ben) {\textbf{Bench}\\ replay robot,\\ every fault};
\node[st=purpleline, right=of rob] (rep) {\textbf{Report}\\ tables, figures\\ from raw files};
\draw[arr] (rec) -- (tr); \draw[arr] (tr) -- (rob); \draw[arr] (tr.east) -- ++(2mm,0) |- (off.west);
\draw[arr] (tr.east) -- ++(2mm,0) |- (ben.west);
\draw[arr] (rob) -- (rep); \draw[arr] (off.east) -- ++(2mm,0) |- (rep.west); \draw[arr] (ben.east) -- ++(2mm,0) |- (rep.west);
\node[below=4mm of rec, font=\scriptsize, text=ink2, align=center] {datasets on the Hub\\ \texttt{recorded.yaml} in git};
\node[below=4mm of tr, font=\scriptsize, text=ink2, align=center] {models on the Hub\\ (best and final checkpoint)};
\node[below=2mm of ben, font=\scriptsize, text=ink2, align=center] {sessions in the\\ results store};
\end{tikzpicture}

  \caption{The path of the project. Datasets are recorded and trimmed on the laptop, policies are trained on the GPU
           servers, and every trained policy is evaluated in three ways: offline, on the robot, and on the replay
           bench. The report is rebuilt from the raw result files.}
  \label{fig:pipeline}
\end{figure}

\begin{table}[htbp]
  \centering
  \small
  \begin{tabularx}{\textwidth}{l X X X}
    \toprule
    & Offline & Robot & Bench \\
    \midrule
    Question & did the policy learn the movements; how fast and large is it? & does the robot do the task under
      each condition? & what does each fault do to the timing and the motion? \\
    Arm moves & no & yes & no (replay robot) \\
    Operator & no & yes: sets the layout, judges success & no \\
    Network & none & the real one (WiFi or cable) & simulated, as measured on the robot \\
    Policy input & held-out recorded frames & the live cameras & held-out episodes, played as if live \\
    Metrics & M9, M10 & M1 and M2--M8, M11, M12 & M2--M8, M11, M12 \\
    \bottomrule
  \end{tabularx}
  \caption{The three kinds of evaluation.}\label{tab:evalkinds}
\end{table}

\paragraph{Offline.}\label{sec:offline-design} Without a robot, each trained policy predicts a chunk from every tenth
frame of every held-out episode, and the first 15 predicted actions (0.5~s, which every policy's chunk covers) are
compared with the 15 recorded actions that follow (M9). Flow and diffusion policies use one sample with a fixed seed.
The same run times 200 requests after 10 warm-up requests and records the peak GPU memory (M10). It runs for the served
checkpoint and the last one, right after each training.

\paragraph{Robot.} The real system end to end: the robot computer (the Raspberry Pi in the official trials) drives the
arm and sends observations over the network (WiFi, or a cable for E4) to a policy server; the operator sets the layout, starts each trial and judges success.

\paragraph{Bench.}\label{sec:bench-design} The same pipeline without the arm: the replay robot plays a held-out episode
as its camera images and joint states, and the policy server runs on the same GPU machine. The network to the robot is
simulated on every reply as measured on the robot: 250~ms with up to 20~ms of jitter (the network share of the
turnaround of the pilot robot run). Each run replays the next held-out episode. Nothing moves and the images do not
react to the actions, so the bench has no task success; it measures what every condition does to the timing and the
motion, under each fault type, many times, without operator time. A run whose control loop falls below 27~Hz is not
counted.

\section{Evaluation plans}\label{sec:evalplans}

Each study is one plan file: its models, its conditions, the number of trials, the written success rule, the time
limit, the comparisons the report makes, and which of them are tested. Before the first trial the plan is
\emph{frozen}: every model is pinned to the exact Hub commit it will be served at. The model servers serve only that
commit, under a name that contains it, so a session cannot reach any other version of a model.

\paragraph{Rounds.} Trials run in rounds. A round runs every model under every condition once, in a random order
seeded by the plan's seed and the round number, so every machine computes the same order. Slow drifts (light, motor
temperature, the operator's attention) therefore affect all conditions alike. One layout of the objects is used for all
trials of a round, and the layouts rotate over the rounds as in the recording.

\paragraph{Labels.} During a trial the operator presses \texttt{s} when the task is done by the written rule (which
also stops the clock for the time to success) or \texttt{f} when it will not succeed; the time limit (30~s) is a
failure. When the control system itself ends the run after the first action (a terminal fault), the trial is a failure
labelled \texttt{fault}, whatever key is pressed afterwards: the label is read from the run's own record, and the key
is kept separately. A fault before the first action is a setup problem: the trial is not counted and runs again in a
later round. A trial with a setup mistake (an object off its mark, a cable caught) is \emph{voided} with a reason; it is
kept and listed, never deleted, and runs again later. Voiding is never used for a policy's failure.

\paragraph{Growth.}\label{sec:growth} A plan with trials may only grow. The number of trials may go up: the next rounds
continue where the plan stood. Models may be added: freezing then pins only the added ones, the frozen ones never
change, and the next rounds run the new models until they have as many trials as the others. Everything else (a
condition, the task, the success rule, the time limit, a model's version) needs a new plan, which the session refuses
to mix with the old one. Stopping at the end of a round keeps every condition at the same number of trials.

\diagram{evalflow}{The evaluation workflow, from a plan to a report.}{fig:evalflow}

\section{Records and the results store}\label{sec:records}

Every trial stores its raw files: one row per control tick (the queue, stalls, the latency estimate, the commanded and
the measured joints, and the policy's own planned action for the step), one row per chunk (send and receive times,
turnaround, server timings, the estimate at sending, the disagreement at the switch), the injected faults, the
scheduler's recovery actions, the trial's configuration (including its fault shift), its label and key, the log, and a
video from each camera. Raw files are written once and never edited. A session is one sitting on one robot computer
with one server; it is pushed to the results store when it ends, and before a session starts the store is pulled, so
the next round accounts for every machine's trials. Reports, tables and figures are always rebuilt from the raw files,
so a change of metrics is one rebuild, and every number in this report can be traced to its trials. Test runs go to a
separate scratch folder that no report reads.

\section{Metrics}\label{sec:metrics}

Table~\ref{tab:metrics} summarises the metrics (version~3); Appendix~\ref{app:metrics} gives the full definitions.
Every metric is first reduced to one value per trial and then summarised over the trials (median and interquartile
range), so long trials do not weigh more than short ones. Times are measured on one machine's monotonic clock.

\begin{table}[htbp]
  \centering
  \small
  \begin{tabularx}{\textwidth}{l l X}
    \toprule
    & Metric & Definition \\
    \midrule
    M1 & Success rate & Successful trials over valid trials, with a Wilson 95\% interval \\
    M1f & Ended by a fault & Trials that the system ended after the first action \\
    M1t & Time to success & From the first action to the success key \\
    M2 & Stall fraction & Share of control ticks after the first action without a new action \\
    M3 & Stalls & Number of stall events; the longest stall in seconds \\
    M4 & Turnaround & From sending an observation to receiving its chunk (median and 95th percentile) \\
    M5 & Estimate error & The scheduler's latency estimate minus the measured turnaround \\
    M6 & Inference time & The server's reported time per chunk; $\mathrm{M4}-\mathrm{M6}$ is the network share \\
    M7 & Jerk & Root mean square of the third time derivative of the commanded arm joints, while the arm moves \\
    M8 & Chunk switch, largest step & Disagreement between chunks at a switch; the largest change between two ticks \\
    M9 & Held-out error & Offline: mean absolute error of the first 15 predicted actions on held-out episodes \\
    M10 & Inference cost & Offline: inference time over 200 requests, and peak GPU memory \\
    M11 & Recovery time & Per fault: from its end until the arm moves again, or \enquote{not recovered} \\
    M12 & Deviation from the plan & Mean distance between the commanded position and the policy's own action \\
    \bottomrule
  \end{tabularx}
  \caption{Metrics (version 3).}\label{tab:metrics}
\end{table}

With $T$ the ticks after the first action and $s_t=1$ when tick $t$ has no new action, the stall fraction is
$\mathrm{M2} = \frac{1}{T}\sum_t s_t$. With $q_{j,t}$ the commanded position of arm joint $j$ on the 30~Hz tick grid
(held through stalls) and $\mathcal{M}$ the jerk samples whose four ticks include at least one executed action, the
jerk is
\begin{equation}
  \mathrm{M7} = \sqrt{\frac{1}{5\,|\mathcal{M}|} \sum_{j=1}^{5} \sum_{t\in\mathcal{M}} \Bigl( f^3\,\bigl(q_{j,t+3}
  - 3q_{j,t+2} + 3q_{j,t+1} - q_{j,t}\bigr) \Bigr)^{2}}, \qquad f = 30~\mathrm{Hz},
\end{equation}
in \jerkunit. Counting only the samples while the arm moves matters: holding still has zero jerk, so over the whole run
a system that stalls more would look smoother. The stop and the restart around a stall still count.

\paragraph{Recovery (M11).} For each timed fault, its end is known from the plan and the trial's shift (a spike ends
when its delayed reply is released). The recovery time is the time from the fault's end until the arm moves again:
the end of the first stall that is still going at that moment or starts within 1~s after it (a reply lost near the end
of a window shows only when the queue runs dry), and 0 when there is none. A run in which the arm never moves again,
or which the system ended before the fault did, is \emph{not recovered}. A run that the operator's key or success ended
first is not counted for that fault.

\paragraph{Deviation from the plan (M12).} Smoothing changes what the arm is sent; M12 measures how far it is from what the
policy asked for. On every executed tick it is the mean over the five arm joints of the distance between the commanded
position and the newest accepted chunk's raw action for that step, before any smoothing. It is 0 without smoothing, and
grows with a filter's lag.

\diagram{metricsdef}{The motion metrics, defined on one trial.}{fig:metricsdef}

\section{Statistics}\label{sec:stats}

Success rates are reported with Wilson 95\% intervals~\cite{wilson1927}. Differences in success between two conditions
of the same policy use Fisher's exact test~\cite{fisher1935}, and differences in the other metrics the two-sided
Mann--Whitney $U$ test~\cite{mann1947} on the per-trial values; in a test of M11 a run that never moved again ranks
after every recovered run. Comparisons are made only within one plan (same sessions, interleaved trials) and only for
the same policy.

With dozens of conditions and metrics, a report can contain thousands of comparisons, and at a 5\% level many would
look significant by chance. Each plan therefore names its comparisons in advance, and a few of them as \emph{primary}:
only those are tested, with the Holm correction~\cite{holm1979} over all primary tests of the plan, and the report also
counts in how many policies the condition is better. Every other comparison is shown as exploratory, with its effect
and an uncorrected $p$-value.

\section{Choosing the smoothing stack}\label{sec:stackrule}

The stack is adaptive playback and the resume ramp, plus one seam method (RTC or temporal ensembling) and one jitter
method (zero-lag smoothing or the Butterworth filter). The two choices are made from the bench by a rule fixed before
the bench ran, and applied by a script. For each policy and dataset, each candidate is compared with \drtc{} alone in
its own plan. A candidate is ruled out for that pair if it adds more than 1\degree{} of median deviation from the plan
(M12), or raises the median stall fraction (M2) or any median recovery time (M11) by more than 10\% (at least 0.5
percentage points and 0.05~s), or recovers in fewer runs. Of the candidates left, the one with the lower jerk relative to
its reference wins the pair; within 1\%, the lower largest step. The candidate winning more pairs is chosen. The guard
on M12 matters because the bench has no task success: without it, the rule would favour whatever smooths hardest,
including by lagging behind the policy, which on the robot can cost a grasp. Policies without RTC use temporal
ensembling as the seam method.

\diagram{stackrule}{The rule that chooses the smoothing stack.}{fig:stackrule}

\section{Summary}
This chapter presents the evaluation framework, combining offline evaluation on held-out episodes, robot trials, and a replay bench without the arm under faults. Each study uses a frozen plan pinning models to Hub commits, running seeded rounds, and growing only by adding rounds or models. Trials save immutable tick, chunk, fault, and video records in a results store from which all reports, tables, and figures are rebuilt. Metrics M1--M12 capture success, latency, stalls, jerk, chunk switches, offline error, recovery, and plan deviation, tested within plans using Fisher's exact test, Mann--Whitney $U$ tests, and the Holm correction. Finally, a fixed script selects the smoothing stack on the bench against \drtc{} alone by gating on deviation, stalls, and recovery before choosing lower jerk.
